\begin{papercard}{alfawzan2022}
\cf{Stages} \stagetag{A} \stagetag{D} \stagetag{I}. An app ``scoping review and content analysis'' of women's mHealth apps on both stores, with a statistically derived popularity cutoff and a yes/no/N-A checklist applied by two reviewers.

\cf{Aim} Assess the privacy policies, data sharing and security policies of popular women's mHealth apps on the App Store and Google Play \pc{alfawzan2022}{1}{The purpose of this scoping review and content analysis was to}. Women's health is defined broadly (sexual/reproductive, physical, mental), not only fertility \pc{alfawzan2022}{2}{and consists of several dimensions}.

\cf{Design and protocol} Follows Arksey and O'Malley's scoping framework, adapted to an app review. A protocol written in November 2020 fixed the search, screening, selection and analysis steps \pc{alfawzan2022}{3}{The scoping review followed the methodology introduced by}.

\cf{Search} Keywords from the literature. Several combinations were tested, starting broad, to find a string with a broad yet specific yield. The search deliberately minimised false negatives at the cost of false positives \pc{alfawzan2022}{4}{The primary database search aimed to}. The string was (woman OR women OR feminine OR female) OR (health OR medical OR medicine), as given in Textbox 1 \pc{alfawzan2022}{4}{Search strings used in the database search}. The search ran in the commercial app-intelligence database 42matters, not in the stores, with four filters: text match in description/developer/title; IAB category ``medical health''; genre medical or health \& fitness; exact word match \pc{alfawzan2022}{4}{The search procedure made use of the}. Run by the first author in January 2021, yielding 136 apps \pc{alfawzan2022}{4}{yielded a total of 136 apps}. The text later says 163, an internal inconsistency worth noticing \pc{alfawzan2022}{4}{Of the 163 apps identified in the search}. Fields recorded per app are listed in Table 1: name, description, downloads (in orders of magnitude, log10), mean rating, rating count (log10), keywords, developer, languages, market status, website, IAB category \pc{alfawzan2022}{4}{Information collected from identified apps}.

\cf{Popularity cutoff (distinctive)} Downloads and rating counts differed strongly between \emph{medical} and \emph{health \& fitness} apps, and many medical apps had no ratings. A single cutoff would be unfair, so separate cutoffs were set per store category \pc{alfawzan2022}{5}{we concluded that the 2 categories needed separate cutoff values}. Rationale: low downloads mean low popularity, low rating counts mean high rating variance, and downloads and ratings correlate \pc{alfawzan2022}{5}{The main reasons for choosing the cutoff values were that}. Method: apps were ordered by log rating count, rating-value variance was computed in sliding windows (bins of 5 to 25), and the point where variance dropped was found visually. For health \& fitness this was bin 21, giving $\geq$720 ratings (16 of 43 apps). For medical apps the same variance value gave $\geq$100 ratings (18 apps) \pc{alfawzan2022}{5}{the rating value variance per bin was calculated for each bin size}. Sensitivity checks inspected the rating counts of the next excluded apps \pc{alfawzan2022}{5}{revealed that the apps not chosen in the sequence had 540 and}. 34 apps passed.

\cf{Exclusions} English availability, functional in January--March 2021, and a real women's-health \emph{service}, not information only (articles, magazines) or games \pc{alfawzan2022}{5}{The app must provide a service related to women}. Removed: four hospital-bound consulting apps, four not actually in English despite the store claim, one not updated for iPhone 11 and crashing on Android, one off-topic, one withdrawn. 23 apps remained \pc{alfawzan2022}{5}{Using these criteria, 4 medical consulting apps were excluded}. A PRISMA flowchart documents the flow \pc{alfawzan2022}{6}{flowchart outlining the search and selection procedure}.

\cf{Assessment set-up} Two reviewers assessed every app independently. Android apps were run in the LDPlayer 4 emulator (Android 7.1.2) and iOS apps on an iPhone 11 (iOS 14.4), between February and March 2021 \pc{alfawzan2022}{6}{assessed independently by 2 reviewers}. Store differences were noted: the App Store shows the policy before download, Google Play did not (at the time), and PEGI age labels were inconsistent with policy age limits \pc{alfawzan2022}{6}{One of the main differences between the Apple App Store and}.

\cf{Instrument} Criteria were taken from two prior app security/privacy studies (Bondaronek et al.\ 2018; Grindrod et al.\ 2016), described as GDPR-aligned, in five domains: privacy policy, data gathering, data sharing, security, transparency \pc{alfawzan2022}{6}{The assessment questions were}. Answers were yes/no, with ``not applicable'' for vague statements or unanswered questions, plus a free-text comment per item that feeds the narrative results \pc{alfawzan2022}{6}{Our assessment was based on yes and no answers}. Example items: policy available in-app or before download; short-form notice; other languages; region-specific sections; contact information; explicit consent at account creation; deletion by request; behaviour and location tracking; third-party, research and legal sharing; security explained; transparency (Table 3) \pc{alfawzan2022}{7}{Is there a short-form notice (in plain English) highlighting key data practices}.

\cf{Reliability} Two independent reviewers, but \textbf{no agreement statistic and no reconciliation procedure are reported}. Only the final counts are given.

\cf{Key findings} 16/23 had an in-app policy, three had none anywhere, and one required at least four clicks via the website \pc{alfawzan2022}{10}{at least four clicks were required to find it}. 11 of 23 apps offered multiple languages but the policy only in English \pc{alfawzan2022}{10}{provided their services in}. Consent: 11 required none. One ``pseudoconsent'' asked users to accept a non-existent policy \pc{alfawzan2022}{11}{We considered this to be a pseudoconsent}. Data were sometimes collected before consent on the welcome page \pc{alfawzan2022}{11}{The welcome pages of 3 other apps}.

\cf{Use for this thesis} Strong precedent for (a) documenting the search platform, date, string and filters, (b) a counted exclusion flow with a PRISMA-style diagram, (c) a \emph{reasoned} popularity threshold rather than an arbitrary one, (d) treating ``policy absent'' and ``N/A'' explicitly, and (e) checking the language of the policy against the app's languages, which is directly relevant for a German-first rule. Limits: third-party database search instead of the stores themselves, category-specific cutoffs that the 13 October meeting rejected (quotas), an unreported agreement statistic and a count inconsistency (136 vs.\ 163).
\end{papercard}
